\documentclass[12pt]{article}
\usepackage[utf8]{inputenc}
\usepackage[T1]{fontenc}
\usepackage[margin=1in]{geometry}
\usepackage{setspace}
\usepackage{xcolor}
\usepackage[colorlinks=true, urlcolor=blue!60!black, linkcolor=blue!60!black]{hyperref}
\setlength{\parindent}{2em}
\setlength{\parskip}{1.25em}
\renewcommand{\baselinestretch}{1.0}
\begin{document}
\begin{center}
{\LARGE \textit{United States v.\ Brown}, 381 U.S.\ 437 (1965)\par}
\vspace{0.8em}
{\large Case Brief --- Four-Page Edition\par}
\end{center}
\vspace{1em}

\section{Memory Jogger}

A longtime Communist Party member elected to a longshoremen's union board was prosecuted under a law barring Party members from union office. The Court held 5--4 that the law was a bill of attainder: Congress itself singled out a named group for punishment instead of enacting a general rule for courts to apply (\href{https://supreme.justia.com/cases/federal/us/381/437/}{Justia}).

\section{Detailed Case Facts}

The Labor-Management Reporting and Disclosure Act of 1959 sought to curb union corruption and protect commerce from ``political strikes.'' Its \S~504 replaced Taft-Hartley's non-Communist affidavit requirement with a crime: any person who was a member of the Communist Party, or had been one in the preceding five years, could not willfully serve as an officer, executive board member, or non-clerical employee of a union, on pain of a \$10,000 fine, a year in prison, or both. Archie Brown, a San Francisco longshoreman, had been an open, avowed Party member for more than twenty-five years. Members of Local 10 of the International Longshoremen's and Warehousemen's Union elected him to their Executive Board in 1959, 1960, and 1961.

\section{Procedural History}

A grand jury indicted Brown in 1961 for willfully violating \S~504. The government never alleged that he had advocated illegal conduct or instigated a political strike, and the trial court excluded evidence that the board had never called a strike. A jury convicted him, and he was sentenced to six months. The en banc Ninth Circuit reversed on First and Fifth Amendment grounds. The Supreme Court affirmed on the alternative ground that \S~504 is a bill of attainder, and did not reach the other issues.

\section{Judicial Votes}

\textbf{5--4}, affirmed (June 7, 1965).
\begin{itemize}
\item \textbf{Majority:} Chief Justice Warren, joined by Justices Black, Douglas, Brennan, and Goldberg.
\item \textbf{Concurrences:} None.
\item \textbf{Dissent:} Justice White, joined by Justices Clark, Harlan, and Stewart.
\end{itemize}

\section{Holding}

Section 504 is an unconstitutional bill of attainder under Article~I, \S~9, cl.~3. The Clause forbids Congress to single out an easily ascertainable group---here, members of the Communist Party---and inflict punishment on them without a judicial trial, even when its aim is preventive (forestalling political strikes) rather than retributive.

\section{Analysis of Opinions}

\textbf{Majority (Warren): history and structure.} The Clause was not a narrow ban on English death sentences but a separation-of-powers safeguard, drawing on Madison, Hamilton, and Marshall in \textit{Fletcher v.\ Peck}. \textit{Cummings}, \textit{Garland}, and \textit{United States v.\ Lovett} confirmed that disqualification from a vocation or office is attainder whether targets are named individually or as a group.

\textbf{Majority (Warren): specification, not rulemaking.} Congress may bar those likely to incite political strikes, but must state general criteria for courts to apply rather than declare that Party members have the feared traits. The Subversive Activities Control Act, by contrast, regulated defined activities with a hearing and review. The Banking Act's bar on underwriters serving as bank directors was different: it did not target a political group, rested on general human psychology, and had an objective standard. Party membership is no neutral ``shorthand'' for dangerousness.

\textbf{Majority (Warren): prevention is punishment.} \textit{American Communications Ass'n v.\ Douds} had called the affidavit rule preventive; Warren repudiated that premise, since punishment serves preventive as well as retributive aims, and many historical attainders against Tories were preventive.

\textbf{White, joined by Clark, Harlan, and Stewart (dissenting).} The majority replaced the multi-factor punishment test of \textit{Flemming v.\ Nestor} and \textit{Mendoza-Martinez} with a per se rule against specificity, threatening conflict-of-interest, felon-disqualification, and security laws and effectively overruling \textit{Douds}. Legislatures regulate specific groups all the time, and \S~504 rationally served the nonpunitive goal of preventing political strikes.

\textbf{Later developments.} \textit{Nixon v.\ Administrator of General Services} (1977) upheld a statute naming President Nixon as a ``legitimate class of one'' and set historical, functional, and motivational tests for punishment; \textit{Selective Service v.\ MPIRG} (1984) applied them.

\section{Examples: Future Applications}

\begin{enumerate}
\item \textbf{Same-Side: Named Protest Group.} A felony bar on members of a named activist coalition serving on utility boards, with a three-year lookback, is an attainder despite its grid-security rationale.
\item \textbf{Same-Side: Salary Rider.} A rider cutting off pay for three named employees accused of disloyalty follows the \textit{Lovett} pattern \textit{Brown} relied on.
\item \textbf{Opposite-Side: Individualized Hearing.} A bar on anyone found, after a hearing with judicial review, to have conspired to sabotage infrastructure is a general rule.
\item \textbf{Opposite-Side: Foreign Vendor.} A procurement ban naming a foreign-adversary-controlled vendor is likely nonpunitive security regulation under \textit{Nixon}.
\item \textbf{Fence-Sitter: Defunding a Named Nonprofit.} Barring a named nonprofit from federal grants looks like attainder under \textit{Brown}, but loss of discretionary funding is not a traditional penalty under \textit{Nixon}.
\end{enumerate}

\section{Critique}

\textbf{Instability after \textit{Nixon}.} \textit{Brown} treats preventive disqualification of an ascertainable group as punishment, yet \textit{Nixon} and later decisions treat preventive national-security targeting as regulation. \textit{Brown} may survive mainly as a rule against criminalizing political-party status. \textbf{First Amendment shadow.} Warren's ``shorthand'' analysis leans on associational-freedom cases such as \textit{Aptheker} and \textit{Schware}, so the structural Attainder Clause may be doing First Amendment work while sidestepping scrutiny tiers.

\section{Key Quotations}

\begin{quote}
``The best available evidence, the writings of the architects of our constitutional system, indicates that the Bill of Attainder Clause was intended not as a narrow, technical (and therefore soon to be outmoded) prohibition, but rather as an implementation of the separation of powers, a general safeguard against legislative exercise of the judicial function, or, more simply, trial by legislature.''
\end{quote}
\noindent --- Chief Justice Warren, for the Court, 381 U.S.\ at 442.

\begin{quote}
``It would be archaic to limit the definition of `punishment' to `retribution.' Punishment serves several purposes; retributive, rehabilitative, deterrent---and preventive. One of the reasons society imprisons those convicted of crimes is to keep them from inflicting future harm, but that does not make imprisonment any the less punishment.''
\end{quote}
\noindent --- Chief Justice Warren, for the Court, 381 U.S.\ at 458.

\begin{quote}
``We do not hold today that Congress cannot weed dangerous persons out of the labor movement, any more than the Court held in \textit{Lovett} that subversives must be permitted to hold sensitive government positions. Rather, we make again the point made in \textit{Lovett}: that Congress must accomplish such results by rules of general applicability. It cannot specify the people upon whom the sanction it prescribes is to be levied. Under our Constitution, Congress possesses full legislative authority, but the task of adjudication must be left to other tribunals.''
\end{quote}
\noindent --- Chief Justice Warren, for the Court, 381 U.S.\ at 461.

\end{document}
